%%%% ABSTRACT
%%
%% Versão do resumo para idioma de divulgação internacional.

\begin{abstractutfpr}%% Ambiente abstractutfpr
Collaborative software development, especially in the digital game industry, relies heavily on large binary files such as three-dimensional models, textures, and audio assets. Traditional version control systems, which are tools designed to track and manage file changes over time, handle these files inefficiently by storing complete copies on every modification, even when only small portions of the content have changed. Existing solutions partially address this issue but still treat binary files as indivisible units, without examining their internal structure. This work proposes a versioning approach that stores only the parts of a three-dimensional binary file that were effectively modified, leveraging its internal structure to identify and isolate each data segment. To detect changes, the SHA-256 function is applied to each segment, generating a unique digital signature that enables precise version comparison. Differential storage is performed through a modern file system, such as Btrfs, which upon recording a modification preserves previous data without overwriting it, sharing unchanged blocks across versions. Experiments compared the proposed method against the conventional approach of storing large files in code repositories, across progressive modification scenarios applied to a three-dimensional model. Initial results indicate expressive reductions in storage consumption, reaching savings of up to 98\% per version in low-change scenarios, and approximately 39\% in total accumulated storage across multiple versions.
\end{abstractutfpr}
